% problem.tex
%
% Purpose:
%   Define the scalar Helmholtz transmission model and the two main workflows:
%   forced lead scattering and homogeneous eigenvalue search.
%
% Main algorithm:
%   The section states the continuous equations and the matrix forms
%   \mathcal A x=b and \mathcal A x=0 used later.
%
% Based on:
%   notes/theory/scat_formulation.md and notes/theory/scat_formulation2.md.
%
% Main changes:
%   The problem statement is condensed into a paper-style formulation.
%
% Numerical goal:
%   Fix the interpretation of the unknowns and right-hand sides before the
%   discretized trace system is introduced.

\section{Transmission and Scattering Problems}

\subsection{Scalar Helmholtz transmission model}

In each homogeneous material region the field satisfies a scalar Helmholtz
equation.  In the exterior background and inside an inclusion we write
\[
  (\Delta+k_e^2)u_e=0,
  \qquad
  (\Delta+k_i^2)u_i=0 .
\]
For a dielectric inclusion with relative index \(n\), one often has
\(k_i=nk_e\), but the formulation only uses the exterior and interior
wavenumbers supplied to the boundary integral assembly.

Across a smooth material interface \(\Sigma_j=\partial\Omega_j\), the TM-like
transmission conditions used in the current formulation are
\[
  u_e=u_i,
  \qquad
  \partial_\nu u_e=\partial_\nu u_i .
\]
Here \(\nu\) is the unit normal pointing out of the inclusion.  A TE variant
would replace the Neumann continuity by a weighted normal derivative condition,
but that is a change in the boundary condition row rather than in the
trace-subspace idea.

The transverse quasiperiodicity is
\[
  u(x,y+d)=\ee^{\ii\beta d}u(x,y),
\]
where \(d\) is the transverse period and \(\beta\) is the transverse
quasimomentum.

\subsection{Forced lead-scattering workflow}

In a forced scattering problem, the data are incoming Bloch traces from one or
both half-leads.  The field is decomposed into prescribed incoming data and an
unknown scattered part that is outgoing or admissible in each half-lead.  The
discretized system has the form
\[
  \Atr(k,\beta)x=b_{\mathrm{in}} .
\]
The right-hand side is not an arbitrary wall trace.  It must be an element of
the incoming trace subspace selected from the periodic half-lead:
\[
  \begin{bmatrix}
  D_{\mathrm{in,tot}}^\pm\\
  N_{\mathrm{in,tot}}^\pm
  \end{bmatrix}
  =
  \begin{bmatrix}
  \Din^\pm\\
  \Nin^\pm
  \end{bmatrix}p_{\mathrm{in}}^\pm .
\]
For a one-sided negative-lead incidence one sets \(p_{\mathrm{in}}^+=0\); for a
one-sided positive-lead incidence one sets \(p_{\mathrm{in}}^-=0\).

\subsection{Homogeneous eigenvalue workflow}

For a defect-mode or transmission-eigenvalue-type calculation, the same assembly
is used without forcing:
\[
  \Atr(k,\beta)x=0 .
\]
Numerically, one scans for parameter values where
\[
  \sigmin(\Atr(k,\beta))\approx 0 .
\]
This homogeneous problem should be interpreted with care.  Away from resonances,
a forced exterior problem is expected to be unique after an appropriate
radiation or limiting absorption condition is imposed.  In contrast, the
eigenvalue workflow is precisely a search for nontrivial admissible null vectors.
Thus the same matrix structure is used for two different purposes: solving
well-posed forced problems and locating lossless or nearly lossless homogeneous
states.
